% GENERATED FILE. Edit canonical lessons, then run npm run generate:books.
% canonical-source-hash: fnv1a64:a4883c7cc599efd2
% canonical-lessons: TA-C241-cataranamana, TA-C241-apurvamana, TA-C241-camamana, TA-C241-karatumuratana, TA-C241-vetuvetuppana

\chapter{Ordinary, Rare, Equal, Rough, Lukewarm}
\label{ch:ta-ordinary-rare-equal-rough-lukewarm}
\hlchaptermodality{241}

\begin{chapteropening}
\textbf{By the end of this chapter:} \emph{I can say ordinary, rare, equal, rough and lukewarm.}

Complete the last lesson of chapter 241: I can say ordinary, rare, equal, rough and lukewarm.
\end{chapteropening}

\section[\texorpdfstring{\ta{சாதாரணமான} (cātāraṇamāṉa)}{சாதாரணமான (cātāraṇamāṉa)}]{\ta{சாதாரணமான} (cātāraṇamāṉa) — ordinary, normal}

\label{lesson:TA-C241-cataranamana}



\begin{quote}
Before the new one: say the Tamil for sharp, then the Tamil for wonderful.
\end{quote}

\subsection*{You'll want to know: \ta{சாதாரணமான}}
\textbf{\ta{சாதாரணமான}} — \emph{cātāraṇamāṉa} — \textquotedblleft{}ordinary, normal\textquotedblright{}.

\begin{cousinweb}
\textbf{\ta{சாதாரணமான}} (\emph{cātāraṇamāṉa}) is \emph{cātāraṇam}, a well-established loan from Sanskrit \emph{sādhāraṇa}, \textquotedblleft{}common, general\textquotedblright{}, with \emph{-āṉa}.
\end{cousinweb}

\begin{grammarlens}[title={What you've built}]
One more word for this chapter.
\end{grammarlens}

\subsection*{Your turn}
\begin{itemize}
\raggedright
  \item \emph{Say it:} \emph{cātāraṇamāṉa}
  \item \emph{Say it:} \emph{cātāraṇamāṉa}, once more
  \item \emph{Say it:} say \emph{cātāraṇamāṉa}
  \item \emph{Recall:} say the Tamil for sharp, then the Tamil for wonderful, then say \emph{cātāraṇamāṉa} again
\end{itemize}

\begin{tcolorbox}[breakable,colback=teal!4,colframe=teal!35!black,title={Before you move on}]
What does \ta{சாதாரணமான} mean? (\textquotedblleft{}Ordinary, normal\textquotedblright{}.) Say it once more.
\end{tcolorbox}
\section[\texorpdfstring{\ta{அபூர்வமான} (apūrvamāṉa)}{அபூர்வமான (apūrvamāṉa)}]{\ta{அபூர்வமான} (apūrvamāṉa) — rare}

\label{lesson:TA-C241-apurvamana}



\begin{quote}
Before the new one: say the Tamil for wonderful, then the Tamil for ordinary.
\end{quote}

\subsection*{You'll want to know: \ta{அபூர்வமான}}
\textbf{\ta{அபூர்வமான}} — \emph{apūrvamāṉa} — \textquotedblleft{}rare\textquotedblright{}.

\begin{cousinweb}
\textbf{\ta{அபூர்வமான}} (\emph{apūrvamāṉa}) is a loan from Sanskrit \emph{apūrva}, \textquotedblleft{}unprecedented\textquotedblright{} (CDIAL 494). \emph{Apūrvam} takes \emph{-āṉa}.
\end{cousinweb}

\begin{grammarlens}[title={What you've built}]
One more word for this chapter.
\end{grammarlens}

\subsection*{Your turn}
\begin{itemize}
\raggedright
  \item \emph{Say it:} \emph{apūrvamāṉa}
  \item \emph{Say it:} \emph{apūrvamāṉa}, once more
  \item \emph{Say it:} say \emph{apūrvamāṉa}
  \item \emph{Recall:} say the Tamil for wonderful, then the Tamil for ordinary, then say \emph{apūrvamāṉa} again
\end{itemize}

\begin{tcolorbox}[breakable,colback=teal!4,colframe=teal!35!black,title={Before you move on}]
What does \ta{அபூர்வமான} mean? (\textquotedblleft{}Rare\textquotedblright{}.) Say it once more.
\end{tcolorbox}
\section[\texorpdfstring{\ta{சமமான} (camamāṉa)}{சமமான (camamāṉa)}]{\ta{சமமான} (camamāṉa) — equal}

\label{lesson:TA-C241-camamana}



\begin{quote}
Before the new one: say the Tamil for ordinary, then the Tamil for rare.
\end{quote}

\subsection*{You'll want to know: \ta{சமமான}}
\textbf{\ta{சமமான}} — \emph{camamāṉa} — \textquotedblleft{}equal\textquotedblright{}.

\begin{cousinweb}
\textbf{\ta{சமமான}} (\emph{camamāṉa}) is a loan from Sanskrit \emph{sama}, \textquotedblleft{}equal, alike\textquotedblright{} (CDIAL 13173). \emph{Camam} takes \emph{-āṉa}.
\end{cousinweb}

\begin{grammarlens}[title={What you've built}]
One more word for this chapter.
\end{grammarlens}

\subsection*{Your turn}
\begin{itemize}
\raggedright
  \item \emph{Say it:} \emph{camamāṉa}
  \item \emph{Say it:} \emph{camamāṉa}, once more
  \item \emph{Say it:} say \emph{camamāṉa}
  \item \emph{Recall:} say the Tamil for ordinary, then the Tamil for rare, then say \emph{camamāṉa} again
\end{itemize}

\begin{tcolorbox}[breakable,colback=teal!4,colframe=teal!35!black,title={Before you move on}]
What does \ta{சமமான} mean? (\textquotedblleft{}Equal\textquotedblright{}.) Say it once more.
\end{tcolorbox}
\section[\texorpdfstring{\ta{கரடுமுரடான} (karaṭumuraṭāṉa)}{கரடுமுரடான (karaṭumuraṭāṉa)}]{\ta{கரடுமுரடான} (karaṭumuraṭāṉa) — rough, uneven}

\label{lesson:TA-C241-karatumuratana}



\begin{quote}
Before the new one: say the Tamil for rare, then the Tamil for equal.
\end{quote}

\subsection*{You'll want to know: \ta{கரடுமுரடான}}
\textbf{\ta{கரடுமுரடான}} — \emph{karaṭumuraṭāṉa} — \textquotedblleft{}rough, uneven\textquotedblright{}.

\begin{cousinweb}
\textbf{\ta{கரடுமுரடான}} (\emph{karaṭumuraṭāṉa}) is native Dravidian. DEDR 1265 lists Kannada \emph{karaḍu}, \textquotedblleft{}rough, uneven\textquotedblright{}. \emph{Karaṭu}, \textquotedblleft{}roughness\textquotedblright{}, is doubled with a rhyming partner, \emph{muraṭu}.
\end{cousinweb}

\begin{grammarlens}[title={What you've built}]
One more word for this chapter.
\end{grammarlens}

\subsection*{Your turn}
\begin{itemize}
\raggedright
  \item \emph{Say it:} \emph{karaṭumuraṭāṉa}
  \item \emph{Say it:} \emph{karaṭumuraṭāṉa}, once more
  \item \emph{Say it:} say \emph{karaṭumuraṭāṉa}
  \item \emph{Recall:} say the Tamil for rare, then the Tamil for equal, then say \emph{karaṭumuraṭāṉa} again
\end{itemize}

\begin{tcolorbox}[breakable,colback=teal!4,colframe=teal!35!black,title={Before you move on}]
What does \ta{கரடுமுரடான} mean? (\textquotedblleft{}Rough, uneven\textquotedblright{}.) Say it once more.
\end{tcolorbox}
\section[\texorpdfstring{\ta{வெதுவெதுப்பான} (vetuvetuppāṉa)}{வெதுவெதுப்பான (vetuvetuppāṉa)}]{\ta{வெதுவெதுப்பான} (vetuvetuppāṉa) — lukewarm}

\label{lesson:TA-C241-vetuvetuppana}



\begin{quote}
Before the new one: say the Tamil for equal, then the Tamil for rough.
\end{quote}

\subsection*{You'll want to know: \ta{வெதுவெதுப்பான}}
\textbf{\ta{வெதுவெதுப்பான}} — \emph{vetuvetuppāṉa} — \textquotedblleft{}lukewarm\textquotedblright{}.

\begin{cousinweb}
\textbf{\ta{வெதுவெதுப்பான}} (\emph{vetuvetuppāṉa}) is native Dravidian (DEDR 5517). The entry lists \emph{vetumpu}, \textquotedblleft{}to become warm\textquotedblright{}, under \emph{vē}, \textquotedblleft{}to burn, to be hot\textquotedblright{}; the doubled \emph{vetuvetu-} makes it \textquotedblleft{}mildly warm\textquotedblright{}.
\end{cousinweb}

\begin{grammarlens}[title={What you've built}]
One more word for this chapter.
\end{grammarlens}

\subsection*{Your turn}
\begin{itemize}
\raggedright
  \item \emph{Say it:} \emph{vetuvetuppāṉa}
  \item \emph{Say it:} \emph{vetuvetuppāṉa}, once more
  \item \emph{Say it:} say \emph{vetuvetuppāṉa}
  \item \emph{Recall:} say the Tamil for equal, then the Tamil for rough, then say \emph{vetuvetuppāṉa} again
\end{itemize}

\begin{tcolorbox}[breakable,colback=teal!4,colframe=teal!35!black,title={Before you move on}]
What does \ta{வெதுவெதுப்பான} mean? (\textquotedblleft{}Lukewarm\textquotedblright{}.) Say it once more.
\end{tcolorbox}
